\begin{proof}
Because the 2D evolution operator factorises as
\(U_{k \times m} = U_k \otimes U_m\) following the tensor-product
decomposition established in Eq.~(\ref{eq:factorised_unitary}),
achieving \(U_{k \times m}^N = U_k^N \otimes U_m^N = e^{i\Phi} I_{4km}\)
requires simultaneous 1D revivals: \(U_k^N = e^{i\phi_x} I_{2k}\) and
\(U_m^N = e^{i\phi_y} I_{2m}\) \cite{Dukes2014}. By Theorem 1, no 1D
revival exists if either \(k\) or \(m\) falls outside the set
\(\{2, 3, 4, 5, 6, 8, 10\}\), thereby bounding the allowable 2D revival
dimensions strictly to Dukes' finite classification. The additional
condition \(\mathcal{N}_k \cap \mathcal{N}_m \neq \emptyset\) ensures
that the two independent 1D walks revive at the same step count \(N\),
a prerequisite for the tensor-product state to return to its initial
configuration up to a global phase.
\end{proof}
